\documentclass[../../../include/open-logic-section]{subfiles}

\begin{document}
\olfileid{sth}{card-arithmetic}{expotough}

\olsection[ځينې ساده کول]{د اصلي توان ځينې ساده کول}

که څه هم د $\canonord$ تعريف لږ پېچلے و، پايله يې د اصلي جمعې
او ضرب په اړه ګټوره ده، \olref[simp]{cardplustimesmax}۔ خو
نامتناهي توان په دومره نېغه توګه نه ساده کېږي۔ د دې د علت د
روښانولو دپاره له يوې پايلې پيل کوو چې د متناهي حالت يوه اشنا
بڼه غځوي (که څه هم ثبوت يې په لوړه انتزاعي کچه دے):

\begin{prop}\ollabel{simplecardexpo}
$\cardexpo{\cardfont{a}}{\cardfont{b} \cardplus \cardfont{c}} =
\cardexpo{\cardfont{a}}{\cardfont{b}} \cardtimes
\cardexpo{\cardfont{a}}{\cardfont{c}}$ او
$\cardexpo{(\cardexpo{\cardfont{a}}{\cardfont{b}})}{\cardfont{c}} =
\cardexpo{\cardfont{a}}{\cardfont{b} \cardtimes \cardfont{c}}$، د
هر درېيو اصلي عددونو $\cardfont{a}, \cardfont{b}, \cardfont{c}$ دپاره۔
\end{prop}

\begin{proof}
د لومړۍ ادعا دپاره يوه تابع $f \colon
(\cardfont{b}\disjointsum\cardfont{c}) \to \cardfont{a}$ وګورئ۔ اوس
دا «ووېشئ»: $f_\cardfont{b}(\beta) = f(\beta, 0)$ د هر
$\beta \in \cardfont{b}$ دپاره، او $f_\cardfont{c}(\gamma) = f(\gamma, 1)$
د هر $\gamma \in \cardfont{c}$ دپاره تعريف کړئ۔ نقشه $f \mapsto
(f_{\cardfont{b}} \times f_\cardfont{c})$ يوه !!a{bijection}
$\funfromto{\cardfont{b} \disjointsum \cardfont{c}}{\cardfont{a}} \to
(\funfromto{\cardfont{b}}{\cardfont{a}} \times
\funfromto{\cardfont{c}}{\cardfont{a}})$ ده۔

د دويمې ادعا دپاره يوه تابع $f \colon \cardfont{c} \to
(\funfromto{\cardfont{b}}{\cardfont{a}})$ وګورئ؛ نو د هر
$\gamma \in
\cardfont{c}$ دپاره يوه تابع $f(\gamma) \colon \cardfont{b} \to
\cardfont{a}$ لرو۔ اوس
$f^*(\beta, \gamma) = (f(\gamma))(\beta)$ د هر
$\tuple{\beta, \gamma} \in \cardfont{b} \times \cardfont{c}$
دپاره تعريف کړئ۔ نقشه
$f \mapsto f^*$ يوه !!a{bijection}
$\funfromto{\cardfont{c}}{(\funfromto{\cardfont{b}}{\cardfont{a}})}
\to \funfromto{\cardfont{b} \cardtimes \cardfont{c}}{\cardfont{a}}$ ده۔
\end{proof}

اوس زموږ \emph{غوښتنه} دا ده چې د نامتناهي اصلي عددونو په حالت
کښې $\cardexpo{\cardfont{a}}{\cardfont{b}}$ په اسانه محاسبه کړو۔
لاندې پايله په دې لوري يو ښه ګام دے:

\begin{prop}\ollabel{cardexpo2reduct}
که $2 \leq \cardfont{a} \leq \cardfont{b}$ وي او $\cardfont{b}$
نامتناهي وي، نو $\cardexpo{\cardfont{a}}{\cardfont{b}} =
\cardexpo{2}{\cardfont{b}}$۔
\end{prop}

\begin{proof}
\begin{align*}
	\cardexpo{2}{\cardfont{b}} &\leq
	\cardexpo{\cardfont{a}}{\cardfont{b}}\text{، ځکه $2 \leq \cardfont{a}$}\\
	&\leq \cardexpo{(2^\cardfont{a})}{\cardfont{b}}
	\text{، د \olref[opps]{lem:SizePowerset2Exp} له مخې}\\
	&= \cardexpo{2}{\cardfont{a} \cardtimes \cardfont{b}}
	\text{، د \olref{simplecardexpo} له مخې} \\
	&= \cardexpo{2}{\cardfont{b}}
	\text{، د \olref[simp]{cardplustimesmax} له مخې}
\end{align*}
\end{proof}

له دې څخه نور ساده کول په رښتيا تمه نه شو کولے، ځکه
$\cardfont{b} < \cardexpo{2}{\cardfont{b}}$ د
\olref[card-arithmetic][opps]{lem:SizePowerset2Exp} له مخې۔
خو دا موږ ته د $\cardexpo{\cardfont{a}}{\cardfont{b}}$ په اړه
نه وايي چې کله $\cardfont{b} < \cardfont{a}$ وي، څه ووايو۔
البته که $\cardfont{b}$ \emph{متناهي} وي،
لار يې پېژنو۔

\begin{prop}
که $\cardfont{a}$ نامتناهي وي او $n \in \omega$ وي، او نوموړے
طبيعي عدد صفر نه وي، نو $\cardexpo{\cardfont{a}}{n} = \cardfont{a}$۔
\textbf{د سرچينې يادښت (OLSTH-019):} چاپي شرط صفر هم رانغاړي؛
په صفر توان کښې يوازې تشه تابع شته او حاصل يو دے، نو ادعا دلته
يوازې د ناصفر متناهي توان دپاره اخيستل شوې ده۔
\end{prop}

\begin{proof}
د ناصفر متناهي توان دپاره
$\cardexpo{\cardfont{a}}{n} = \cardfont{a} \cardtimes \cardfont{a}
\cardtimes \ldots \cardtimes \cardfont{a} = \cardfont{a}$، د
\olref[simp]{cardplustimesmax} له مخې۔
\end{proof}
\noindent
پر دې سربېره، په ځينو نورو حالتونو کښې هم د
$\cardexpo{\cardfont{a}}{\cardfont{b}}$ اندازه ټاکلے شو:

\begin{prop}
که $2 \leq \cardfont{b} < \cardfont{a} \leq
\cardexpo{2}{\cardfont{b}}$ وي او $\cardfont{b}$ نامتناهي وي، نو
$\cardexpo{\cardfont{a}}{\cardfont{b}} = \cardexpo{2}{\cardfont{b}}$۔
\end{prop}

\begin{proof}
$\cardexpo{2}{\cardfont{b}}\leq \cardexpo{\cardfont{a}}{\cardfont{b}}
\leq \cardexpo{(\cardexpo{2}{\cardfont{b}})}{\cardfont{b}} =
\cardexpo{2}{\cardfont{b}\cardtimes\cardfont{b}} =
\cardexpo{2}{\cardfont{b}}$، لکه په \olref{cardexpo2reduct} کښې
استدلال وشو۔
\end{proof}
\noindent
خو له دې ځايه وروسته مسئلې ډېرې نازکې کېږي۔

\end{document}
